%% ============================================================
%% 分片 25：原文第 9 节「't Hooft And Hecke Operators」（原文 glshort.tex 起于第 6755 行）
%% 由 tools/harvmac2tex.py 从 arXiv 源码（hep-th/0604151v3）自动生成。
%% 只译英文散文；公式、编号、\label、\cite 一律不改。
%% ============================================================

\subsection{仿射 Grassmann 流形}\label{sec:affgrass}

与 't Hooft 算子相关联的 Hecke 修正空间，可以描述为某个称为仿射 Grassmann 流形的无穷维簇中的一个有限维子簇。这一描述使人能够为 Hecke 修正空间构造自然的紧化。本文其实不会真正使用这个框架，但我们在这里解释它，因为它提供了理解某些事情的强有力方式，而且在数学上这个课题通常就被安置在这个框架里。物理学家熟悉仿射 Grassmann 流形，主要是由于它在二维共形场论中的应用 \cite{ref122}。

为清楚起见，我们考虑 $G=U(N)$ 的情形。此时，一个 't Hooft 算子由模去置换的一组整数 $m_1,\ldots,m_N$ 指定。由于问题是局部的，我们可以把曲线 $C$ 固定为 $\Bbb{CP}^1$，并取 $E_-$ 为平凡丛。$E_-$ 在点 $p_0\in C$ 处的 Hecke 修正 $E$，是这样一个丛：它带有在 $p_0$ 之外的一个选定同构 $\sigma:E\cong E_-$；$\sigma$ 可能无法延拓过 $p_0$。

自然地把固定 $m_1,\ldots,m_N$ 的 Hecke 修正空间看作配对 $(E,\sigma)$ 构成的空间的子空间，其中 $E$ 是 $C$ 上的全纯向量丛，$\sigma$ 是它在 $p_0$ 之外的平凡化。下面沿用 \cite{ref122} 第 8 章，回顾这种配对 $(E,\sigma)$ 构成的空间的显式描述。

方便的做法是把 $C=\Bbb{CP}^1$ 看作复 $z$-平面 $\CC$ 的单点紧化，并取 $p_0$ 为点 $z=0$。我们还令 % 式 (9.28)
\setcounter{equation}{27}
\begin{equation}\label{eq:jmj}
U_\infty=\Bbb{CP}^1-\{0\},\quad U_0=\Bbb{CP}^1-\{\infty\} .
\end{equation}
我们在 $U_\infty$ 上给定 $E$ 的一个平凡化 $\sigma$。虽然 $\sigma$ 可能无法延拓到 $U_0$ 上，但我们可以在 $U_0$ 上取 $E$ 的一个平凡化 $\sigma'$。在 $U_0\bigcap U_\infty\simeq\CC^*$ 上，$\sigma$ 与 $\sigma'$ 相差一个 $GL(N,\CC)$ 规范变换 $g(z)$。这是取值于 $GL(N,\CC)$ 的函数，其矩阵元在 $\CC^*$ 上全纯，但在 $0$ 与 $\infty$ 处可以有极点。（总可以选取 $\sigma'$，使 $g(z)$ 的奇点都是极点。）用 ${\cal X}$ 记在 $\CC^*$ 上全纯、并具有到 $\Bbb{CP}^1$ 的亚纯延拓的函数环；于是 $\CC^*$ 上的全纯规范变换群可以与 $GL(N,{\cal X})$ 等同，而 $g(z)$ 是这个群的一个元素。

若用在整个 $U_0$ 上全纯的规范变换 $h_0(z)$ 改变 $\sigma'$，则 $g(z)$ 被替换为 $g(z)h_0^{-1}(z)$。$h_0(z)$ 取值于 $GL(N,\cO)$，即矩阵元为 $U_0\simeq \CC$ 上全纯函数的可逆矩阵群。$\Bbb{CP}^1$ 上秩为 $N$ 的、在 $p_0$ 之外带有固定全纯平凡化 $\sigma$ 的全纯向量丛之集合，同构于商 $GL(N,{\cal X})/GL(N,\cO)$。后一空间称为 $GL(N)$ 的仿射 Grassmann 流形 ${\rm Gr}_N$。它还有一个名字叫圈 Grassmann 流形（loop Grassmannian），因为它同构 \cite{ref122}\ 于 $U(N)$ 中基于给定基点的圈构成的空间。

\begin{yizhu}[易错点]
按：$g(z)$ 不是丛的自同构，而是两个平凡化 $\sigma$、$\sigma'$ 之间的过渡函数。换 $U_0$ 上的平凡化 $\sigma'$ 对应右乘 $h_0^{-1}$；而换 $E_+$ 的生成截面（即换 $f_j$）对应左乘 $GL(N,\cO)$ 的元素。Hecke 修正空间取的是 $GL(N,\cO)$ 在 $g$ 上的左作用轨道，右乘的部分已被商 $GL(N,{\cal X})/GL(N,\cO)$ 吸收掉；把左右作用弄反会得到不同的轨道分解。
\end{yizhu}

仿射 Grassmann 流形的定义允许一些变体。例如，由于 ${\cal X}$ 中任何函数都可以展开为 Laurent 级数，可以把 ${\cal X}$ 嵌入形式 Laurent 级数环 $\CC((z))$。类似地，可以把全纯函数环 $\cO$ 嵌入形式幂级数环 $\CC[[z]]$。可以证明商 $GL(N,\CC((z)))/GL(N,\CC[[z]])$ 同构于仿射 Grassmann 流形 ${\rm Gr}_N$。直观上，之所以如此，是因为给定 $GL(N,{\cal X})$ 的任意两个元素，只需研究它们 Laurent 展开中有限多项，就能判断它们是否属于同一个 $GL(N,\cO)$ 轨道，因此这些 Laurent 级数是否有非零的收敛区域并不重要。

现在我们必须指出 ${\rm Gr}_N$ 中哪些点可以由平凡向量丛经权为 $^L\negthinspace w=(m_1,\ldots,m_N)$ 的 Hecke 修正得到。以这个权对平凡向量丛作 Hecke 修正，就是选取平凡向量丛在 $U_0$ 上的一个由 $N$ 个线性无关的全纯截面 $f_1,f_2,\dots,f_N$ 给出的平凡化，并声明 $E_+$ 在 $U_0$ 上由截面 $s_1,\ldots,s_N$ 生成，而这些截面限制到 $U_\infty\bigcap U_0$ 上时由下式给出 % 式 (9.29)
\setcounter{equation}{28}
\begin{equation}\label{eq:mcop}
z^{-m_j} f_j,\quad j=1,\ldots,N.
\end{equation}
这里我们用到：在 $U_\infty$ 上给定了 $E_+$ 与平凡向量丛之间的一个同构。于是，若 $f_j=(f_j^1,f_j^2,\dots,f_j^N)$，则对应于 $E_+$ 的矩阵 $g(z)$ 由下式给出 % 式 (9.30)
\setcounter{equation}{29}
\begin{equation}\label{eq:xcoxo}
g^i_j(z)=z^{-m_j}
f^i_j(z).
\end{equation}
最简单的选择是 $f^i_j(z)=\delta^i_j$；所有其他选择都可以从这一个通过用 $GL(N,\cO)$ 的元素左乘 $g(z)$ 得到。由此得出：Hecke 修正空间 ${{\EUY}}({}^L\negthinspace w;p_0)$ 就是点 % 式 (9.31)
\setcounter{equation}{30}
\begin{equation}\label{eq:gspec}
g^i_j(z)=z^{-m_j}
\delta^i_j
\end{equation}
在仿射 Grassmann 流形中、群 $GL(N,\cO)$ 左作用下的轨道。可以把 ${\cal X}$ 处处换成 $\CC((z))$、把 $\cO$ 换成 $\CC[[z]]$，得到等价的结果。注意 \eqref{eq:gspec}\ 给出的 $g(z)$ 是 $\CC^*$ 到 $GL(N,\CC)$ 的同态。事实上，它是同态 % 式 (9.32)
\setcounter{equation}{31}
\begin{equation}\label{eq:xxconon}
\rho: U(1)\to U(N),\quad \rho:
e^{i\alpha}\to \left(\begin{matrix}e^{i m_1\alpha} & & \ldots & \\ &
e^{i m_2\alpha} & \ldots &  \\  & \vdots & \ldots & \\ & & &
e^{i m_N\alpha}\\\end{matrix}\right)
\end{equation}
的复化，而这个同态进入 't Hooft 算子 $T(m_1,\ldots,m_N)$ 的定义。

我们已把 $GL(N,\cO)$ 在 ${\rm Gr}_N$ 中的一个轨道与每一组整数 $m_1,\ldots,m_N$ 对应起来。结果表明，${\rm Gr}_N$ 中所有 $GL(N,\cO)$ 轨道都是这样得到的 \cite{ref122} 。于是 ${\rm Gr}_N$ 被空间 ${{\EUY}}(m_1,\ldots,m_N)$ 分层，这些空间称为 Schubert 胞腔。等价地，对平凡向量丛施加 Hecke 修正，可以得到 $\Bbb{CP}^1$ 上任意的全纯向量丛，并带有 $U_\infty$ 上任意的平凡化。

\begin{yizhu}[易错点]
按：轨道其实只依赖于整数集 $m_1,\ldots,m_N$ 的置换类。置换矩阵的矩阵元为常数，因而属于 $GL(N,\cO)$；重排 $m_i$ 只是把轨道共轭一下，得到的是同一个轨道族。这与 't Hooft 算子的权本就定义在置换之下（对应 Weyl 群的作用）是一致的，所以原文一开始就说 $m_1,\ldots,m_N$ 是「模去置换」给定的。
\end{yizhu}

第~\ref{sec:spaceheck}~节 讨论的 Hecke 修正空间的例子提示这些空间总是有限维的。这一点一般地可以如下证明。设 ${\rm Gr}_N(i)$ 是 ${\rm Gr}_N$ 中由下述条件定义的子集：$g(z)$ 在 $z=0$ 处的极点阶数不高于 $i$。簇 ${\rm Gr}_N(i)$ 是有限维的、$GL(N,\cO)$-不变的且紧的，而 ${\rm Gr}_N$ 中任何点都属于某个 $i$ 的 ${\rm Gr}_N(i)$。这意味着 $GL(N,\cO)$ 的任何轨道都属于某个 ${\rm Gr}_N(i)$，从而是有限维的。

此时不妨提一下：无穷维空间 ${\rm Gr}_N$ 还有另一种分层，其层由整数集合 $k_1,\ldots,k_N$ 标记。为定义这种分层，我们回忆 Grothendieck 的一个定理：$\Bbb{CP}^1$ 上任何秩为 $N$ 的全纯向量丛都同构于 % 式 (9.33)
\setcounter{equation}{32}
\begin{equation}\label{eq:Groth}
\oplus_{i=1}^N
\cO(p_0)^{k_i}
\end{equation}
其中 $k_1,\ldots,k_N$ 是某些整数。我们把 ${\rm Gr}_N$ 中对应于 $k_1,\ldots,k_N$ 的层定义为满足下述条件的配对 $(E,\sigma)$ 之集合：$E$ 同构于 \eqref{eq:Groth}。换句话说，层就是固定 $E$ 而变动平凡化 $\sigma$ 得到的。显然，这里的层都是无穷维的，与 ${{\EUY}}(m_1,\ldots,m_N)$ 给出的分层形成对照。${\rm Gr}_N$ 的这两种分层之间的关系在 \cite{ref122} 中研究过。

\begin{yizhu}[含糊点拨]
按：${\rm Gr}_N$ 上有两套容易混淆的分层。一套按全纯向量丛的拓扑型 $\oplus_{i=1}^N\cO(p_0)^{k_i}$ 分（Grothendieck 定理），层是无穷维的：固定 $E$、让平凡化 $\sigma$ 变动。另一套按 Schubert 胞腔 ${{\EUY}}(m_1,\ldots,m_N)$ 分，层是有限维的：在 Hecke 修正的框架里，$E_-$ 与它在 $U_\infty$ 上的平凡化都固定，变动的只是 $E_+$。前者标记「丛是什么」，后者标记「Hecke 修正的权」。
\end{yizhu}

如我们所见，Hecke 修正空间或 Schubert 胞腔一般是非紧的。紧化它们的自然方式是：对充分大的 $i$，考虑相应轨道在 ${\rm Gr}_N(i)$ 中的闭包。对给定的 $\rho: U(1)\to U(N)$，轨道 $C_\rho$ 的闭包一般是一个奇异簇，它是 $C_\rho$ 与有限多个更低维轨道的并。它称为 Schubert 循环。正如第~\ref{sec:spaceheck}~节 所讨论的，$C_\rho$ 闭包的结构反映了不同 $\rho$ 但同一拓扑型的 't Hooft 算子之间的「混合」。我们还将在第~\ref{sec:monbub}~节 用单极子冒泡来描述这一过程。

\begin{yizhu}[物理图像]
按：Schubert 循环（轨道闭包）就是 Hecke 修正空间的自然紧化。闭包里出现的更低维轨道正是「混合」的几何表现：它们与原来的轨道有相同的拓扑型（例如相同的 $\oplus\cO(p_0)^{k_i}$），却对应不同的 $\rho$；在闭包的边界上，不同权的 't Hooft 算子无法再区分。物理图像是单极子冒泡——单极子缩成一点并并入已有的单极子，从而把一个权换成另一个权（第~\ref{sec:monbub}~节会详细展开）。
\end{yizhu}

${\rm Gr}$ 的构造以及把 Hecke 修正空间描述为 ${\rm Gr}$ 的子簇这一做法，可以推广到其他规范群。对任何紧单 Lie 群 $G$，记其复化为 $G_{\Bbb{C}}$，把 ${\rm Gr}_G$ 定义为商 % 式 (9.34)
\setcounter{equation}{33}
\begin{equation}\label{eq:hxonoh}
{\rm Gr}_G=G_{\Bbb{C}}({\cal
X})/G_{\Bbb{C}}(\cO).
\end{equation}
't Hooft 算子由同态 $\rho: U(1)\to G$ 的共轭类参数化，而 $\rho$ 可以解析延拓为同态 $\rho_\CC: \CC^*\to G_{\Bbb{C}}$。显然，$\rho_\CC$ 定义了 ${\rm Gr}_G$ 上的一个点，而该点在 $G_{\Bbb{C}}(\cO)$ 左作用下的轨道只依赖于 $\rho$ 的共轭类。这个轨道就是 Schubert 胞腔或 Hecke 修正空间 ${{\EUY}}(\rho)$。

\begin{yizhu}[易错点]
按：这里起决定作用的是复化群 $G_{\Bbb{C}}$，不是紧群 $G$ 本身：${\rm Gr}_G$ 定义中的商是 $G_{\Bbb{C}}({\cal X})/G_{\Bbb{C}}(\cO)$。$\rho_\CC:\CC^*\to G_{\Bbb{C}}$ 是 $\rho:U(1)\to G$ 的解析延拓（$U(1)$ 上的同态由其在单位元处的值决定，故可延拓）。对 $G=U(N)$，把 $\rho$ 对角化为 $\mathrm{diag}(e^{im_1\alpha},\ldots,e^{im_N\alpha})$，其共轭类就由 $(m_1,\ldots,m_N)$ 的置换类给出，回到前面的整数集。
\end{yizhu}
